\chapter{Tổng quan đề tài}
\label{Chapter1}

\section{Bối cảnh và vấn đề nghiên cứu}

Trong những năm gần đây, các ứng dụng web được phát triển với tốc độ ngày càng nhanh nhằm đáp ứng nhu cầu thay đổi liên tục của doanh nghiệp và người dùng. Các mô hình phát triển phần mềm hiện đại như Agile và DevOps rút ngắn chu kỳ phát hành, khiến mỗi hệ thống có thể được cập nhật nhiều phiên bản trong thời gian ngắn. Điều này đặt ra yêu cầu kiểm thử thường xuyên để đảm bảo các chức năng hiện có vẫn hoạt động chính xác sau mỗi lần thay đổi.

Trong quy trình kiểm thử phần mềm, việc xây dựng test case, chuẩn bị test data và viết automation test script thường chiếm nhiều thời gian và chi phí. Đối với các dự án có số lượng chức năng lớn hoặc phát hành nhiều phiên bản, tester phải lặp lại quá trình phân tích yêu cầu, thiết kế test case và hiện thực test script cho từng tính năng. Khi hệ thống thay đổi, các test script hiện có cũng cần được bảo trì hoặc chỉnh sửa, làm tăng đáng kể chi phí kiểm thử.

Sự phát triển của các mô hình ngôn ngữ lớn (Large Language Models - LLM) mở ra khả năng hỗ trợ tự động sinh test case và test data từ yêu cầu bằng ngôn ngữ tự nhiên. Tuy nhiên, nếu mỗi lần kiểm thử hoặc mỗi lần phát hành phiên bản mới đều phụ thuộc hoàn toàn vào LLM hoặc AI Agent để sinh lại toàn bộ kịch bản kiểm thử thì chi phí sử dụng mô hình và thời gian suy luận có thể tăng, đồng thời quy trình thực thi sẽ phụ thuộc nhiều hơn vào kết quả suy luận của agent.

Do đó, bài toán đặt ra không chỉ là tự động hóa việc sinh test case hay test script, mà còn cần xây dựng một quy trình cho phép tester tận dụng khả năng của AI để giảm khối lượng công việc ban đầu, đồng thời tạo ra các test script có khả năng tái sử dụng qua nhiều phiên bản phần mềm nhằm giảm chi phí bảo trì và kiểm thử hồi quy.

Từ thực tế đó, đề tài tập trung nghiên cứu và xây dựng một hệ thống hỗ trợ kiểm thử tự động website theo hướng kết hợp mô hình ngôn ngữ lớn với AI Agent. Trong quy trình này, AI hỗ trợ sinh test case và test script, còn tester đóng vai trò kiểm tra, lựa chọn và xác nhận kết quả trước khi đưa vào sử dụng. Các test script sau khi được tạo và xác nhận sẽ được lưu trữ để tái sử dụng trong các lần kiểm thử tiếp theo, hạn chế việc phải sinh lại bằng AI ở mỗi phiên bản.

\section{Mục tiêu đề tài}

Mục tiêu tổng quát của đề tài là xây dựng một hệ thống hỗ trợ kiểm thử tự động website nhằm giảm thời gian xây dựng test case, test script và test data, đồng thời giảm chi phí phát triển và bảo trì automation testing trong quá trình phát triển phần mềm.

Các mục tiêu cụ thể bao gồm:

\begin{itemize}
    \item Xây dựng nền tảng web hỗ trợ quản lý project kiểm thử và website mục tiêu.
    
    \item Hỗ trợ tester sinh test case từ yêu cầu hoặc mô tả bằng ngôn ngữ tự nhiên thông qua mô hình ngôn ngữ lớn.

    \item Hỗ trợ tester lựa chọn, chỉnh sửa và xác nhận các test case trước khi thực thi.

    \item Sử dụng Browser-use Agent kết hợp Playwright để tự động sinh automation test script từ các test case đã được xác nhận.

    \item Cho phép tester chạy thử, đánh giá và hoàn thiện phiên bản đầu tiên của test script.

    \item Lưu trữ test script để tái sử dụng trong các lần kiểm thử hồi quy của những phiên bản tiếp theo, giảm sự phụ thuộc vào mô hình ngôn ngữ lớn.

    \item Quản lý test data, test suite, object repository, kết quả chạy kiểm thử và các bằng chứng thực thi phục vụ phân tích và truy vết.
\end{itemize}

\section{Phạm vi đề tài}

Đề tài tập trung vào hỗ trợ kiểm thử giao diện người dùng (UI Testing) đối với các ứng dụng web chạy trên trình duyệt máy tính để bàn.

Đối tượng sử dụng chính của hệ thống là tester hoặc kỹ sư kiểm thử, những người cần xây dựng và duy trì các automation test script trong quá trình phát triển phần mềm.

Trong phạm vi đề tài, hệ thống hỗ trợ quy trình gồm các bước: sinh test case từ yêu cầu bằng ngôn ngữ tự nhiên, lựa chọn và xác nhận test case, sinh automation test script bằng AI Agent, thực thi kiểm thử, lưu trữ kết quả và tái sử dụng test script cho các phiên bản tiếp theo.

Hệ thống hướng đến môi trường staging hoặc môi trường kiểm thử có kiểm soát. Các trường hợp như CAPTCHA, OTP, xác thực đa yếu tố, thanh toán thực hoặc các cơ chế chống bot phức tạp không thuộc phạm vi nghiên cứu.

Đề tài cũng không tập trung vào kiểm thử hiệu năng, kiểm thử bảo mật, kiểm thử API độc lập, kiểm thử ứng dụng di động hoặc kiểm thử tải.

\section{Giải pháp đề xuất}

Đề tài đề xuất một quy trình kiểm thử tự động theo hướng kết hợp giữa mô hình ngôn ngữ lớn (LLM), AI Agent và cơ chế tái sử dụng automation test script.

Đầu tiên, tester nhập yêu cầu hoặc mục tiêu kiểm thử bằng ngôn ngữ tự nhiên. Mô hình ngôn ngữ lớn phân tích yêu cầu và sinh ra danh sách các test case cùng test data đề xuất. Tester xem xét, lựa chọn và xác nhận các test case phù hợp trước khi thực hiện bước tiếp theo.

Sau khi được xác nhận, Browser-use Agent sử dụng các test case để tự động thao tác trên trình duyệt và sinh automation test script thông qua Playwright. Tester có thể chạy thử các script được sinh ra, kiểm tra kết quả và điều chỉnh nếu cần. Phiên bản script sau khi được xác nhận sẽ được lưu trữ như phiên bản chuẩn phục vụ kiểm thử.

Trong kiểm thử hồi quy, hệ thống cho phép tester chủ động lựa chọn chạy lại bằng replay script đối với các kịch bản đã ổn định, thay vì sử dụng agent để suy luận lại toàn bộ luồng. Khi các thay đổi của website vượt quá khả năng tái sử dụng, tester có thể sử dụng AI để tạo mới hoặc cập nhật test script.

Cách tiếp cận này tạo điều kiện giảm số lần gọi mô hình trong các lần chạy hồi quy, đồng thời tăng khả năng tái sử dụng và bảo trì các kịch bản kiểm thử trong quá trình phát triển phần mềm.

\section{Kết quả đạt được}

Sau quá trình thực hiện, đề tài đã xây dựng được một hệ thống web hỗ trợ quy trình tạo và quản lý automation testing theo hướng kết hợp giữa mô hình ngôn ngữ lớn và AI Agent.

Các kết quả chính đạt được gồm:

\begin{itemize}
    \item Xây dựng hệ thống quản lý project, test case, test script, test suite, test data và kết quả kiểm thử.

    \item Tích hợp mô hình ngôn ngữ lớn để hỗ trợ sinh test case và test data từ yêu cầu bằng ngôn ngữ tự nhiên.

    \item Cho phép tester xem xét, lựa chọn và xác nhận test case trước khi sinh automation script.

    \item Tích hợp Browser-use Agent kết hợp Playwright để sinh automation test script và thực thi kiểm thử trên trình duyệt.

    \item Hỗ trợ lưu trữ, quản lý và tái sử dụng test script trong các lần kiểm thử hồi quy.

    \item Ghi nhận log, screenshot, evidence và báo cáo phục vụ truy vết và phân tích kết quả kiểm thử.

    \item Ghi nhận số token sử dụng trong quá trình sinh test case, sinh test data và thực thi bằng agent; từ đó làm cơ sở ước tính chi phí LLM theo mô hình được cấu hình.

    \item Xây dựng quy trình hỗ trợ giảm thời gian tạo automation test, tăng khả năng tái sử dụng replay script và giảm mức phụ thuộc vào LLM trong các lần kiểm thử hồi quy.
\end{itemize}

Nhìn chung, hệ thống đã hiện thực được quy trình hỗ trợ tester từ giai đoạn phân tích yêu cầu, sinh test case, sinh automation test script, xác nhận kết quả đến tái sử dụng các script trong các phiên bản tiếp theo. Qua đó góp phần giảm thời gian xây dựng kiểm thử tự động, nâng cao hiệu quả kiểm thử hồi quy và tạo cơ sở định lượng mức sử dụng LLM thông qua số token được ghi nhận.

\section{Cấu trúc báo cáo}

Báo cáo gồm bảy chương.

Chương 1 trình bày bối cảnh nghiên cứu, mục tiêu, phạm vi, giải pháp đề xuất và kết quả đạt được của đề tài.

Chương 2 trình bày cơ sở lý thuyết và các công nghệ liên quan, bao gồm kiểm thử phần mềm, kiểm thử tự động, mô hình ngôn ngữ lớn, AI Agent và các công cụ hỗ trợ browser automation.

Chương 3 phân tích yêu cầu hệ thống, các nhóm người dùng, yêu cầu chức năng, yêu cầu phi chức năng và các use case chính.

Chương 4 trình bày kiến trúc và thiết kế hệ thống, bao gồm kiến trúc tổng thể, thiết kế dữ liệu và các luồng xử lý chính.

Chương 5 trình bày quá trình hiện thực hệ thống, bao gồm frontend, backend, cơ sở dữ liệu, quy trình sinh test case, sinh test script và thực thi kiểm thử.

Chương 6 trình bày thực nghiệm và đánh giá hệ thống thông qua các kịch bản kiểm thử trên website mẫu, đồng thời đánh giá hiệu quả về thời gian tạo test case, khả năng tái sử dụng replay script và mức sử dụng token LLM trong kiểm thử hồi quy.

Chương 7 trình bày kết luận, những hạn chế của đề tài và các hướng phát triển trong tương lai.
